% TEMPLATE-MILC-v3.tex - plantilla maestra UNICA de las guias MILC v3.
%
% La usan las cuatro familias de guias, cada una con su driver:
%   · Grados 6-11          -> scripts/build-guias-g11.py
%   · Bebras               -> scripts/build-guias-bebras.py
%   · Semillero            -> scripts/build-guias-semillero.py
%   · Territorio interior  -> scripts/build-guias-territorio-interior.py
%
% Antes habia tres copias casi identicas (~400 lineas iguales, 6 distintas):
% cada mejora habia que aplicarla tres veces, y en la practica no se hacia
% (el motor de recursos vivia solo en dos de ellas). Lo que varia entre
% familias esta parametrizado en 6 placeholders que cada driver llena
% (se nombran aqui SIN su sintaxis real: el builder reemplaza por texto plano,
% tambien dentro de los comentarios, y un valor multilinea rompe el preambulo):
%
%   HEADER_IZQ        encabezado izquierdo de pagina
%   HEADER_DER        encabezado derecho de pagina
%   PORTADA_ETIQUETA  rotulo sobre el numero grande (GRADO/MOMENTO/MODULO)
%   PORTADA_SERIE     linea de serie de la portada
%   PORTADA_ID        identificador de la guia en la portada
%   PORTADA_CIRCULO   el circulito de grado (°), o vacio si no aplica
%   PDF_TITULO        titulo en texto plano (sin \textbf ni comillas TeX) para
%                     los metadatos del PDF (/Title); lo produce lib_xelatex.texto_pdf
%
% Composicion (auditoria 2026-09-03): las bandas de estacion piden espacio
% (\Needspace*) y prohiben el salto justo despues (\nopagebreak), asi no quedan
% huerfanas al pie; las cajas largas (softbox, drawbox, infoband, puentebox) son
% `breakable`, asi una caja de media pagina ya no arrastra todo a la siguiente
% dejando la anterior al 40 %. Todo texto sobre mostaza o verde va en milcNegro
% (blanco sobre #E5B400 da 1,93:1 y sobre #58A923 2,95:1; ninguno pasa WCAG AA).
% El resto de claves las documenta content/guias/_SCHEMA.md. El script valida
% que no queden placeholders sin reemplazar antes de escribir el archivo final.

% Etiquetado PDF (latex-lab): /Lang, estructura y texto alternativo de las
% figuras, para lectores de pantalla. Debe ir ANTES de \documentclass.
\DocumentMetadata{lang=es-CO,tagging=on}
\documentclass[11pt,letterpaper]{article}
% headheight=24pt: la cabecera derecha lleva dos lineas (identidad de la guia
% y estacion actual). headsep se acorta en la misma medida para que el bloque
% de texto quede exactamente donde estaba con headheight=12pt.
\usepackage[margin=2.0cm,headheight=24pt,headsep=13pt]{geometry}
\usepackage[table]{xcolor}
\usepackage{fontspec}
\usepackage[spanish]{babel}
\usepackage{tikz}
% Librerías TikZ para los diagramas de las guías (bloque `recursos:`):
%  · babel  → babel en español vuelve activos caracteres como «>», y TikZ los usa
%             en sus opciones (>=stealth, ->). Sin esta librería, cualquier
%             diagrama con flechas revienta con «\language@active@arg> has an
%             extra }». Debe ir ANTES de que se dibuje nada.
%  · positioning        → sintaxis «below left=2pt and 3pt».
%  · decorations.pathreplacing → llaves ({brace}) para señalar rangos.
\usetikzlibrary{babel,positioning,decorations.pathreplacing}
\usepackage{graphicx}
% Circuitos: el trabajo de electrónica se hace hoy con micro:bit y sus sensores
% integrados (MakeCode, sin cableado), así que no se necesitan esquemáticos.
% Si algún día entran montajes con protoboard, basta descomentar esta línea:
% los diagramas .tex/.tikz declarados en `recursos:` ya se incrustan solos.
% \usepackage{circuitikz}
\usepackage[most]{tcolorbox}
\usepackage{tabularx}
\usepackage{array}
\usepackage{fancyhdr}
\usepackage{titlesec}
\usepackage[document]{ragged2e}  % bandera a la derecha en todo el cuerpo
\usepackage{enumitem}
\usepackage{microtype}
\usepackage{etoolbox}
\usepackage{needspace}
% hyperref al final del preambulo: metadatos del PDF (titulo, autor, idioma,
% familia), marcadores por estacion (\pdfbookmark) y enlaces sin recuadro.
\usepackage[hidelinks,
  pdftitle={Lo que le debo al lugar: lo que recibiste no lo hiciste tú},
  pdfauthor={PhD. Álvaro Cárdenas Orozco},
  pdfsubject={Territorio interior · Educación socioemocional · Grado 11° · Momento 8 · MILC v3},
  pdflang={es-CO},
  bookmarksopen=true,bookmarksnumbered=false]{hyperref}

% Helvetica Neue no trae la flecha U+2192, asi que cada «→» escrito literalmente
% en el YAML se compilaba a NADA, en silencio (462 flechas en 61 guias: rutas del
% tipo «paleta → tipografia → lockup» salian rotas). Mapeamos el caracter a su
% version matematica, que si tiene glifo. Asi el autor puede seguir escribiendo
% «→» comodamente en el YAML.
\usepackage{newunicodechar}
\newunicodechar{→}{\ensuremath{\rightarrow}}
% Mismo problema con los simbolos de marca y vinetas: el «✓» de «una palomita ✓»
% salia como cuadrito vacio en el PDF de todas las guias que lo usaban.
\usepackage{pifont}
\newunicodechar{✓}{\ding{51}}
\newunicodechar{✗}{\ding{55}}
\newunicodechar{✔}{\ding{52}}
\newunicodechar{•}{\textbullet}
\newunicodechar{≥}{\ensuremath{\geq}}
\newunicodechar{≤}{\ensuremath{\leq}}

\setmainfont{Helvetica Neue}
\setsansfont{Helvetica Neue}
\setlength{\parindent}{0pt}
\setlength{\parskip}{6pt}
% Sin particion de palabras: en 6.º-7.º una palabra cortada por guion al final
% de linea («Comuni-cacion») frena la lectura mucho mas que una linea corta.
\hyphenpenalty=10000
\exhyphenpenalty=10000
\pretolerance=2000
\tolerance=3000
\setlength{\emergencystretch}{3em}
\linespread{1.10}
\renewcommand{\arraystretch}{1.25}
\setlist{leftmargin=5mm,itemsep=1mm,topsep=1mm}
\newcolumntype{Y}{>{\RaggedRight\arraybackslash}X}

\definecolor{milcMagenta}{HTML}{E6007E}
\definecolor{milcVino}{HTML}{5A0038}
\definecolor{milcTurquesa}{HTML}{0093A5}
\definecolor{milcVerde}{HTML}{58A923}
\definecolor{milcMostaza}{HTML}{E5B400}
\definecolor{milcOcre}{HTML}{B86600}
\definecolor{milcNegro}{HTML}{241816}
\definecolor{milcGris}{HTML}{F7F7F4}
\definecolor{milcLinea}{HTML}{E8E2DD}
\definecolor{milcGradePrimary}{HTML}{0D9488}
\definecolor{milcGradeDark}{HTML}{115E59}
\definecolor{milcGradeSoft}{HTML}{CCFBF1}
\definecolor{milcGradeAccent}{HTML}{CFE7FF}
\definecolor{milcGradeText}{HTML}{FFFFFF}
\color{milcNegro}

\pagestyle{fancy}
\fancyhf{}
% Cabecera de dos lineas: identidad de la guia arriba y, a la derecha y en
% gris, la estacion en curso (\markright desde \milcestacion). Antes de la
% primera estacion la segunda linea va vacia; el \strut mantiene la altura
% para que la linea de arriba no baile entre paginas.
\lhead{\small Territorio interior · Educación socioemocional\\[-1pt]{\footnotesize\strut}}
\rhead{\small Grado 11° · Momento 8 · MILC v3\\[-1pt]{\footnotesize\strut\color{milcNegro!65}\rightmark}}
\cfoot{\small\thepage}
\renewcommand{\headrulewidth}{0pt}

% Sobre mostaza (#E5B400) y verde (#58A923) el texto va en milcNegro; sobre el
% resto de la paleta, en blanco. Se usa dentro de fonttitle/fontupper, donde un
% \color posterior manda sobre coltitle.
\newcommand{\milctintasobre}[1]{%
  \ifboolexpr{test{\ifstrequal{#1}{milcMostaza}} or test{\ifstrequal{#1}{milcVerde}}}%
    {\color{milcNegro}}{\color{white}}}

\newtcolorbox{titlebox}[1]{
  enhanced,colback=#1,colframe=#1,arc=7mm,boxrule=0pt,
  left=7mm,right=7mm,top=3mm,bottom=3mm,
  before skip=8pt,after skip=8pt,
  fontupper=\sffamily\bfseries\Large\color{white}
}

\newtcolorbox{softbox}[3]{
  enhanced,breakable,pad at break=3mm,
  colback=#2,colframe=#1,borderline west={4pt}{0pt}{#1},
  arc=2mm,boxrule=.4pt,left=4mm,right=4mm,top=3mm,bottom=3mm,
  before skip=6pt,after skip=8pt,title={#3},coltitle=white,
  colbacktitle=#1,fonttitle=\sffamily\bfseries\milctintasobre{#1},fontupper=\RaggedRight,
  attach boxed title to top left={xshift=3mm,yshift=-2mm},
  boxed title style={arc=2mm, boxrule=0pt}
}

\newtcolorbox{drawbox}[2]{
  enhanced,breakable,pad at break=4mm,
  colback=white,colframe=#1,arc=4mm,boxrule=1pt,
  left=5mm,right=5mm,top=4mm,bottom=4mm,before skip=8pt,after skip=8pt,
  title={#2},coltitle=white,colbacktitle=#1,fonttitle=\sffamily\bfseries\milctintasobre{#1},
  fontupper=\RaggedRight,
  attach boxed title to top left={xshift=4mm,yshift=-2mm},
  boxed title style={arc=3mm, boxrule=0pt}
}

\newtcolorbox{infoband}[1]{
  enhanced,breakable,pad at break=2mm,colback=#1!8,colframe=#1!50,arc=2mm,boxrule=.5pt,
  left=4mm,right=4mm,top=2mm,bottom=2mm,before skip=4pt,after skip=4pt,
  fontupper=\small\RaggedRight
}

% Puente narrativo entre secciones (contrato editorial Conéctate).
% Da continuidad: "Acabas de X. Ahora vas a Y porque Z."
% Visualmente: banda izquierda en color del grado, texto en itálica pequeña.
\newtcolorbox{puentebox}{
  enhanced,breakable,pad at break=2mm,colback=milcGradeSoft,colframe=milcGradeDark,
  borderline west={3pt}{0pt}{milcGradeDark},
  arc=0mm,boxrule=0pt,left=5mm,right=4mm,top=2.5mm,bottom=2.5mm,
  before skip=4pt,after skip=8pt,
  fontupper=\itshape\small\color{milcNegro}\RaggedRight
}
% La flecha va en modo matematico a proposito: Helvetica Neue NO tiene el glifo
% U+2192, asi que \textrightarrow se compilaba a NADA («Missing character: There
% is no → in font Helvetica Neue») y el puente salia sin flecha. En modo
% matematico la flecha viene de la fuente matematica y si se dibuja.
\newcommand{\puente}[1]{%
  \begin{puentebox}%
    \textcolor{milcGradeDark}{$\rightarrow$}\hspace{2mm}#1%
  \end{puentebox}%
}

\newcommand{\linead}[1]{\noindent\color{milcLinea}\rule{#1}{0.5pt}\color{milcNegro}}
\newcommand{\respuesta}[1]{\par\vspace{2mm}\foreach \n in {1,...,#1}{\noindent\color{milcLinea}\rule{\linewidth}{0.5pt}\par\vspace{5mm}}\color{milcNegro}}
\newcommand{\checkbox}{\textcolor{milcGradeDark}{\fbox{\phantom{x}}}\hspace{5pt}}

% ─── Listas aireadas para las guías socioemocionales ────────────────────
% Componer las verificaciones como «\checkbox texto\\» deja el texto que salta
% de línea debajo del cuadrito y apelmaza el bloque. Estos entornos alinean la
% continuación y airean los ítems. Los usa Territorio Interior; cualquier
% familia puede hacerlo.
\newlist{checklista}{itemize}{1}
\setlist[checklista]{
  label=\textcolor{milcGradeDark}{\fbox{\phantom{x}}},
  leftmargin=9mm, labelsep=3.2mm, itemsep=2.4mm,
  topsep=2.5mm, parsep=0pt, partopsep=0pt, before=\RaggedRight
}
\newlist{pasoslista}{enumerate}{1}
\setlist[pasoslista]{
  label=\textcolor{milcGradeDark}{\sffamily\bfseries\arabic*.},
  leftmargin=8mm, labelsep=3mm, itemsep=2.8mm,
  topsep=1mm, parsep=0pt, partopsep=0pt, before=\RaggedRight
}

\newcommand{\phasecard}[4]{
\begin{tcolorbox}[enhanced,colback=#3,colframe=#2,arc=3mm,boxrule=.5pt,height=3.05cm,valign=center,halign=center,left=1.5mm,right=1.5mm,top=2mm,bottom=2mm]
{\sffamily\bfseries\fontsize{22}{24}\selectfont\textcolor{#2}{#1}}\par
{\sffamily\bfseries\small #4}
\end{tcolorbox}
}

% ── Piezas del rediseno 2026 (legibilidad 6.º-11.º) ───────────────────
% Banda de estacion: reemplaza el titlebox macizo. Titulo a la izquierda,
% dato practico (tiempo/modalidad) a la derecha, en una sola linea.
% Nunca huerfana: pide 9 lineas libres (o salta de pagina rellenando con
% \vfill) y prohibe el salto justo despues de la banda. Nueve y no seis:
% la caja partible que sigue necesita ~80pt (titulo adosado, relleno y dos
% lineas) para arrancar en la misma pagina; con menos, tcolorbox la manda
% entera a la siguiente y la banda se queda sola. El penalty 10000 va
% en `after`, ANTES del espacio posterior: un `after skip` (aunque sea 0pt)
% es un \vskip tras la caja, y ese glue es un punto de corte legal que un
% \nopagebreak escrito despues de \end{tcolorbox} ya no cierra. Cada banda
% deja marcador PDF y actualiza la cabecera derecha.
\newcounter{milcestacion}
\newcommand{\milcestacion}[2]{%
\Needspace*{9\baselineskip}%
\stepcounter{milcestacion}%
\pdfbookmark[1]{#2}{milcestacion\themilcestacion}%
\markright{#2}%
\begin{tcolorbox}[enhanced,colback=#1,colframe=#1,arc=2mm,boxrule=0pt,
  left=5mm,right=5mm,top=2.6mm,bottom=2.6mm,before skip=11pt,
  after={\par\nopagebreak\vskip7pt}]
{\sffamily\bfseries\fontsize{15}{18}\selectfont\milctintasobre{#1}#2}
\end{tcolorbox}}

% Tarjeta de paso numerada: sustituye las tablas «Pilar / Aplicacion», que
% eran formato de consulta docente, no de estudiante. El numero va en un
% circulo sobre el borde para que la secuencia se lea de un vistazo.
\newcommand{\milcpaso}[3]{%
\Needspace{4\baselineskip}%
\begin{tcolorbox}[enhanced,colback=white,colframe=milcLinea,arc=1.5mm,boxrule=.9pt,
  left=14mm,right=4mm,top=3mm,bottom=3mm,before skip=4pt,after skip=4pt,
  fontupper=\RaggedRight,
  overlay={\node[anchor=north west,circle,fill=#2,text=white,inner sep=0pt,
    minimum size=7.4mm,font=\sffamily\bfseries\small]
    at ([xshift=3.6mm,yshift=-3.2mm]frame.north west) {#1};}]
#3
\end{tcolorbox}}

% Casilla marcable de tamano util con lapiz (la anterior era \fbox{\phantom{x}},
% de alto variable segun el texto que la seguia).
\newcommand{\milcmarca}{\framebox[3.8mm]{\rule{0pt}{3.4mm}}}

% Fila marcable de lista suelta (checks de la estacion 1).
\newcommand{\milccheckrow}[1]{%
\par\vspace{1.8mm}\noindent
\makebox[6.4mm][l]{\raisebox{-.55mm}{\textcolor{milcNegro}{\milcmarca}}}%
\parbox[t]{\dimexpr\linewidth-6.4mm\relax}{\RaggedRight #1}\par}

% Celda marcable dentro de la rubrica (tabla de autoevaluacion). Misma
% sangria francesa, pero pensada para vivir dentro de una columna X.
\newcommand{\milcopcion}[1]{%
\makebox[5.2mm][l]{\raisebox{-.55mm}{\textcolor{milcNegro}{\milcmarca}}}%
\parbox[t]{\dimexpr\linewidth-5.2mm\relax}{\RaggedRight #1}}

% ── Recursos visuales de la guía ──────────────────────────────────────
% Declarados en el bloque `recursos:` del YAML y emitidos por el builder.
% \guiaFigura   → imagen rasterizada o PDF (foto, carta celeste, montaje).
% \guiaDiagrama → fuente TikZ/circuitikz (.tex) incrustada como vector:
%                 es la vía para circuitos y esquemáticos editables.
% [#1] = texto alternativo (alt) · #2 = ruta relativa al .tex · #3 = pie
% El alt llega al PDF etiquetado (/Alt) y lo lee el lector de pantalla.
\newcommand{\guiaFigura}[3][]{%
\begin{center}
\includegraphics[width=0.88\linewidth,height=0.38\textheight,keepaspectratio,alt={#1}]{#2}\\[1.5mm]
{\footnotesize\itshape\textcolor{milcNegro!75}{#3}}
\end{center}
}

\newcommand{\guiaDiagrama}[2]{%
\begin{center}
\input{#1}\\[1.5mm]
{\footnotesize\itshape\textcolor{milcNegro!75}{#2}}
\end{center}
}

\begin{document}
\thispagestyle{empty}

% ── Portada ───────────────────────────────────────────────────────────
% Antes: un rectangulo de color a pagina completa. Se veia bien en pantalla,
% pero en la fotocopiadora del colegio se comia el toner y salia gris sucio.
% Ahora la banda ocupa el tercio superior y el resto va en blanco.
\begin{tikzpicture}[remember picture,overlay]
  \path[fill=milcGradePrimary]
    (current page.north west) rectangle ([yshift=-10.6cm]current page.north east);

  \node[anchor=north west,text=milcGradeText] at ([xshift=2.0cm,yshift=-1.55cm]current page.north west)
    {\sffamily\bfseries\fontsize{12}{14}\selectfont MOMENTO};
  \node[anchor=north west,text=milcGradeText,opacity=.85] at ([xshift=2.0cm,yshift=-2.35cm]current page.north west)
    {\sffamily\bfseries\fontsize{8.5}{10}\selectfont TERRITORIO INTERIOR · SOCIOEMOCIONAL};
  \node[anchor=north east,text=milcGradeText,opacity=.85,align=right] at ([xshift=-2.0cm,yshift=-1.55cm]current page.north east)
    {\sffamily\bfseries\fontsize{9}{12}\selectfont I.E. Sor María Juliana\\Cartago, Valle del Cauca};

  \node[anchor=west,text=milcGradeText] (gradonum) at ([xshift=1.85cm,yshift=-7.1cm]current page.north west)
    {\sffamily\bfseries\fontsize{112}{112}\selectfont 8};
  % El circulito de grado (°) solo aplica a las guias de grado («11°»); en
  % Bebras y Semillero el numero grande es un momento/modulo, no un grado.
  % Se posiciona relativo al nodo `gradonum`, no en coordenadas absolutas.
  

  \node[anchor=west,text=milcGradeText,text width=11.4cm,align=left] at ([xshift=1.5cm,yshift=0.5cm]gradonum.east)
    {
     {\sffamily\bfseries\fontsize{9}{11}\selectfont\MakeUppercase{Territorio interior · Socioemocional}}\\[3mm]
     {\sffamily\bfseries\fontsize{21}{25}\selectfont\setlength{\baselineskip}{27pt}Lo que le\\[4pt]debo al lugar\\[4pt]esto no lo hiciste tú}\\[3.5mm]
     {\sffamily\bfseries\fontsize{15}{18}\selectfont Grado 11° · Momento 8}
    };
\end{tikzpicture}

\vspace*{9.4cm}
\vfill

{\sffamily\bfseries\fontsize{8}{10}\selectfont\textcolor{milcGradeDark}{LO QUE TE LLEVAS HOY}}

\begin{tcolorbox}[enhanced,colback=white,colframe=milcNegro,arc=2mm,boxrule=1.6pt,
  left=5mm,right=5mm,top=4mm,bottom=4mm,before skip=3pt,after skip=12pt,
  fontupper=\RaggedRight\fontsize{11}{15.5}\selectfont]
Tu inventario de lo heredado: diez cosas que usas y no hiciste, con quién las hizo y quién las mantiene hoy; la que está en peor estado; y un compromiso de devolución concreto, con fecha, hacia algo que vas a seguir usando.
\end{tcolorbox}

\begin{tcolorbox}[enhanced,colback=milcGris,colframe=milcGris,arc=2mm,boxrule=0pt,
  left=5mm,right=5mm,top=3.5mm,bottom=3.5mm,after skip=12pt]
{\sffamily\bfseries\fontsize{8}{10}\selectfont\textcolor{milcNegro!55}{ESTA GUÍA ES DE}}\\[3mm]
\begin{tabularx}{\linewidth}{X p{3.2cm} p{3.2cm}}
\linead{\linewidth} & \linead{3.0cm} & \linead{3.0cm}\\[-1mm]
{\fontsize{8.5}{10}\selectfont\textcolor{milcNegro!55}{Nombre completo}} &
{\fontsize{8.5}{10}\selectfont\textcolor{milcNegro!55}{Grupo}} &
{\fontsize{8.5}{10}\selectfont\textcolor{milcNegro!55}{Fecha}}\\
\end{tabularx}
\end{tcolorbox}

\vfill

\noindent\textcolor{milcLinea}{\rule{\linewidth}{1pt}}\\[2mm]
{\sffamily\bfseries\fontsize{9.5}{12}\selectfont Autor: Dr. Álvaro Cárdenas Orozco}\\[1mm]
{\sffamily\fontsize{8.5}{11}\selectfont\textcolor{milcNegro!55}{Metodología MILC · Apertura ancestral · Territorio y deuda · Triángulo Dussel-Estoicismo-Floridi}}

\newpage

\section*{Apertura: Lo que le debo al lugar: lo que recibiste no lo hiciste tú}

\begin{softbox}{milcGradeDark}{milcGradeSoft}{Saber ancestral}
En lengua nasa yuwe, la tierra se dice \textbf{kiwe}. \emph{Y aparece siempre acompañada:} \textbf{\emph{mai kiwe}}, madre tierra; \textbf{\emph{kwe'sx kiwe}}, nuestra tierra; \textbf{\emph{nasa kiwe}}, la tierra del pueblo nasa. \textbf{Los nasa describen su relación con kiwe como una relación de \emph{hermandad}} ---no de dueño a cosa, sino entre parientes---. \emph{Y eso no es una manera bonita de hablar: organiza cómo se decide.} \textbf{Hay \emph{kiwe thë}, los \emph{sabios de la tierra}, que orientan espiritualmente ese vínculo,} \emph{y las plantas ---\textbf{tasxkwe}--- se entienden como \textbf{personas}, con conocimiento propio.} \textbf{Piensa en la diferencia práctica que hace esa distinción.} \emph{Con una propiedad, uno hace lo que quiere: la usa, la explota, la vende, la deja acabar.} \textbf{Con un pariente, no.} \emph{Con un pariente hay obligaciones que uno no escogió ---llegaron con el parentesco---, hay que responder por él, y no se le trata según convenga.} \textbf{Y aquí está lo que interesa para esta guía:} \emph{una deuda con un pariente \textbf{no se paga y se acaba} ---se sostiene toda la vida---.} \textbf{Eso es lo contrario de una factura, y es de lo que trata este momento.}
\end{softbox}

\begin{softbox}{milcTurquesa}{milcGris}{Saber contemporáneo}
¿Y qué hace que algo de todos dure generaciones en vez de derrumbarse? \textbf{Se estudió, y con mucho cuidado.} \textbf{Elinor Ostrom} y su equipo revisaron \textbf{cientos de casos} en los años ochenta ---bosques, aguas de riego, pesquerías, pastos--- buscando qué reglas eran las mejores para administrar un \textbf{bien común}. \emph{No encontró un conjunto de reglas que sirviera siempre.} \textbf{Encontró otra cosa, publicada en \emph{Governing the Commons} (1990): \textbf{ocho principios de diseño} que caracterizaban a los sistemas de bienes comunes \emph{que habían durado mucho tiempo}.} \textbf{Entre ellos:} \emph{que las reglas las hagan \textbf{quienes usan el recurso}, ajustadas a las necesidades y condiciones de esa comunidad; que quienes vigilan \textbf{respondan ante los usuarios}; y que existan \textbf{mecanismos baratos} para resolver los conflictos.} \textbf{Y ahora el hallazgo que convierte esto en el tema de hoy:} \emph{esos principios \textbf{se cumplían} en los sistemas que habían durado, y \textbf{estaban ausentes} en los que se derrumbaron.} \textbf{Es decir: lo que dura no dura solo.} \emph{Alguien, en cada generación, hizo el trabajo de sostenerlo ---y donde nadie lo hizo, el sistema se cayó---.}
\end{softbox}

\begin{softbox}{milcMostaza}{milcGris}{Pregunta-puente}
\textit{Los nasa tratan la tierra \textbf{como pariente y no como propiedad}, y con un pariente hay obligaciones que uno no escogió. \textbf{¿Qué usas todos los días que no hiciste tú} \ldots \textbf{y quién lo está sosteniendo ahora mismo}?}
\end{softbox}

\begin{softbox}{milcVerde}{milcGris}{Saber y saber hacer}
Al terminar podrás: (1) \textbf{reconocer} cuánto de lo que usas a diario existe porque otros lo hicieron y alguien lo mantiene; (2) \textbf{distinguir} tratar algo como propiedad de tratarlo como pariente, y qué obligaciones distintas produce cada cosa; (3) \textbf{explicar} por qué un bien común dura solo si cada generación hace el trabajo de sostenerlo; (4) \textbf{comprometerte} con una devolución concreta y fechada hacia algo que vas a seguir usando.
\end{softbox}



\small
\begin{tabularx}{\textwidth}{>{\bfseries}p{3.0cm}Y}
\rowcolor{milcGradePrimary!12}Autor & Dr. Álvaro Cárdenas Orozco\\
DBA & Comprende los fenómenos que afectan a su territorio, regula sus emociones ante ellos y acompaña a otros con empatía (transversal 6°--11°, articulado con la política «Escuela, territorio de vida» del MEN, Resolución 006519 de 2025, y la Ley 1620 de 2013)\\
Referentes & Primera ayuda psicológica (OMS, 2011/2012) · Directrices IASC (2007) · USGS y Servicio Geológico Colombiano · Pueblo nasa y la minga (CRIC) · Dussel · Estoicismo · Floridi\\
Producto final & Tu inventario de lo heredado: diez cosas que usas y no hiciste, con quién las hizo y quién las mantiene hoy; la que está en peor estado; y un compromiso de devolución concreto, con fecha, hacia algo que vas a seguir usando.\\
\end{tabularx}
\normalsize

\puente{En el momento 7 viste dónde se decide lo común. \textbf{Hoy miras de dónde salió lo común que ya está}, \emph{y qué pasa si nadie de tu generación lo sostiene}.}

\milcestacion{milcGradeDark}{Ruta MILC: así vamos a aprender}
\begin{tabularx}{\textwidth}{>{\centering\arraybackslash}X>{\centering\arraybackslash}X>{\centering\arraybackslash}X>{\centering\arraybackslash}X}
\phasecard{1}{milcVerde}{milcGris}{Escucha\\[-1mm]\small Observo, escucho y pregunto.} &
\phasecard{2}{milcTurquesa}{milcGris}{Sistematización\\[-1mm]\small Recibo y me hago cargo.} &
\phasecard{3}{milcMagenta}{milcGris}{Praxis\\[-1mm]\small Construyo y aplico.} &
\phasecard{4}{milcVino}{milcGris}{Evaluación\\[-1mm]\small Reflexión liberadora.}
\end{tabularx}


\puente{Primero, un inventario que casi nadie ha hecho: qué usas y no hiciste.}

\milcestacion{milcVerde}{Estación 1 de 3 · Escucha}

\begin{softbox}{milcVerde}{milcGris}{Escena para escuchar}
\textbf{Actividad 1 · IDENTIFICA --- Lo que uso y no hice.} \textbf{Primera parte.} Haz una lista de \textbf{diez cosas que usaste ayer y que no hiciste tú}. \emph{Empieza por lo obvio y sigue por lo invisible:} \textbf{el agua que salió de la llave, la vía por la que llegaste, la luz, el colegio, la lengua en la que estás pensando, una vacuna, el parque, una receta, un oficio que alguien te enseñó, el hospital que existe por si acaso.} \textbf{Segunda parte.} Al frente de cada una escribe, si puedes, \emph{quién la hizo} ---no un nombre necesariamente: «gente que peleó por eso», «la junta del barrio», «mi abuela», «el Estado», «no sé»---. \textbf{Cuenta cuántas quedaron en «no sé».} \textbf{Tercera parte.} Ahora la pregunta que casi nadie se hace: \emph{¿quién la mantiene \textbf{hoy}?} \emph{Alguien barre ese parque, alguien opera ese acueducto, alguien enseña esa lengua, alguien pelea por que ese hospital siga.} \textbf{Y la última:} \emph{¿alguna de esas diez cosas está peor que hace unos años?} \textbf{¿Cuál?}
\end{softbox}

{\sffamily\bfseries\fontsize{8}{10}\selectfont\textcolor{milcNegro!55}{MARCA CUANDO LO HAYAS HECHO}}
\milccheckrow{Listaste diez cosas concretas, incluidas las invisibles.}
\milccheckrow{Anotaste \textbf{quién la hizo}, y contaste cuántas quedaron en «no sé».}
\milccheckrow{Respondiste \textbf{quién la mantiene hoy}.}

\begin{infoband}{milcVerde}


\textbf{En tu cuaderno (Actividad 1 · IDENTIFICA).} Título: «Actividad 1. Lo que uso y no hice». \textbf{Formato:} las diez cosas con quién las hizo, quién las mantiene y cuál está peor. \textbf{Extensión:} media página. \textbf{Sabes que terminaste cuando} sabes \textbf{cuál está en peor estado}.
\end{infoband}

\puente{Ya lo tienes. Ahora, la diferencia entre propiedad y pariente, y qué hace que un bien común dure o se derrumbe.}

\milcestacion{milcTurquesa}{Fase 2 · Sistematización: entender la herencia}

\begin{softbox}{milcTurquesa}{milcGris}{Cuatro claves sobre lo heredado}
Cuatro claves para dejar de tratar lo que hay como si simplemente estuviera ahí.
\end{softbox}

\milcpaso{1}{milcTurquesa}{\textbf{Casi nada de lo que usas lo hiciste tú, y eso no es una acusación.} \emph{Es una descripción:} \textbf{el agua potable, la vía, la escuela pública, la vacuna, el alfabeto, la electricidad, el derecho a estudiar siendo mujer, el derecho a no trabajar a los doce años.} \emph{Nada de eso existía por defecto ---todo lo consiguió alguien, y varias de esas cosas costaron peleas largas---.} \textbf{Y lo que hace este inventario es corregir una ilusión muy común a los diecisiete:} \emph{sentir que uno empieza de cero y que todo lo que tenga será mérito propio}. \textbf{Nadie empieza de cero.} \emph{Se empieza dentro de algo que otros construyeron, y saberlo no quita mérito ---ubica---.}}
\milcpaso{2}{milcTurquesa}{\textbf{Pariente, no propiedad: dos maneras de estar en un lugar.} \emph{Con una \textbf{propiedad}, la relación es de uso:} \textbf{se aprovecha mientras rinda y se deja cuando no.} \emph{Con un \textbf{pariente} ---que es como los nasa nombran su relación con kiwe--- la relación es de \textbf{obligación mutua}:} \textbf{hay deberes que uno no escogió y que llegaron con el vínculo.} \emph{Y fíjate en la consecuencia más práctica de esa diferencia:} \textbf{con una propiedad, irse no rompe nada; con un pariente, irse no cancela la relación.} \emph{Por eso este momento no le pide a nadie que se quede ---eso ya lo decidiste en el momento 5---:} \textbf{le pide que sepa que la relación con el lugar no se acaba porque uno se vaya.} \emph{Mucha gente que se fue de esta región sigue sosteniendo cosas de aquí.}}
\milcpaso{3}{milcTurquesa}{\textbf{Lo común se sostiene o se cae, y está estudiado.} \emph{Ostrom buscó las mejores reglas y no las encontró; encontró \textbf{ocho principios de diseño} presentes en los sistemas de bienes comunes que habían durado} ---entre ellos, \emph{que las reglas las hagan quienes usan el recurso, que quienes vigilan respondan ante los usuarios y que resolver un conflicto sea barato}--- \textbf{y \emph{ausentes} en los que se derrumbaron} (Ostrom, 1990). \emph{Léelo con calma, porque es lo contrario de lo que uno supone:} \textbf{lo que dura no dura por inercia.} \emph{Un acueducto veredal, un salón comunal, un camino, una biblioteca, un río limpio, una tradición ---todos se caen en una generación si nadie hace el trabajo---.} \textbf{Y casi siempre ese trabajo lo hacen pocas personas, sin que se note, hasta que dejan de hacerlo.}}
\milcpaso{4}{milcTurquesa}{\textbf{Devolver no es agradecer, y confundirlos deja todo igual.} \emph{Agradecer es un sentimiento y se agota en quien lo siente.} \textbf{Devolver es una acción y le llega a alguien más.} \emph{Y hay una forma específica de devolución que este momento propone y que se distingue de la caridad:} \textbf{no es dar algo a alguien que está peor ---eso puede ser bueno y es otra cosa---,} \emph{sino \textbf{hacerse cargo de una parte de lo que uno recibió, para que siga existiendo}.} \textbf{Es exactamente la lógica de la minga del grado 8:} \emph{la cuenta no se lleva en dinero sino en turnos, y el que hoy recibió mañana va a la de otro.} \emph{Por eso el compromiso de hoy es hacia algo que vas a \textbf{seguir usando}:} \textbf{devolver dentro del sistema del que uno hace parte, no desde afuera.}}

\begin{drawbox}{milcTurquesa}{Cómo se mira lo que uno recibió}
\begin{pasoslista}
\item \textbf{¿Existía antes de mí?} \emph{Casi todo.}
\item \textbf{¿Quién lo hizo, y le costó algo?} \emph{Muchas cosas costaron peleas que ya nadie recuerda.}
\item \textbf{¿Quién lo mantiene hoy?} \emph{Si la respuesta es «no sé» o «nadie», está en riesgo.}
\item \textbf{¿Lo estoy tratando como propiedad o como pariente?}
\item \textbf{¿Qué parte de esto podría sostener yo?} \emph{Pequeña, concreta, sostenida ---no heroica---.}
\end{pasoslista}
\end{drawbox}

\begin{softbox}{milcMostaza}{milcGris}{Errores comunes y su corrección}
\textbf{Creer que uno empieza de cero:} se empieza dentro de algo que otros construyeron. \textbf{Confundir agradecer con devolver:} el primero se agota en uno. \textbf{Confundir devolver con caridad:} devolver es hacerse cargo dentro del sistema del que uno hace parte. \textbf{Pensar que lo que existe seguirá existiendo:} lo que dura, lo sostiene alguien. \textbf{«Cuando tenga plata devuelvo»:} casi todo lo que sostiene un lugar no cuesta plata sino tiempo. \textbf{Prometer sin fecha:} una deuda sin fecha no se paga. \textbf{Creer que irse cancela el vínculo:} con un pariente, no.
\end{softbox}

\begin{infoband}{milcTurquesa}


\textbf{En tu cuaderno (Actividad 2 · EXPLICA).} Título: «Actividad 2. Propiedad o pariente». \textbf{Formato:} dos columnas con las obligaciones que produce cada relación, y debajo explica qué encontró Ostrom sobre lo que distingue a los sistemas que duraron de los que se derrumbaron. \textbf{Extensión:} media página. \textbf{Sabes que terminaste cuando} puedes explicar por qué \textbf{lo que dura no dura solo}.
\end{infoband}

\puente{Ya sabes lo que recibiste. Es hora de la parte incómoda: qué de eso está en peor estado, y qué vas a hacer tú.}

\milcestacion{milcMagenta}{Fase 3 · Praxis: hacerse cargo}

\begin{softbox}{milcMagenta}{milcGris}{Cómo se devuelve algo que uno recibió}
Un inventario, una constatación incómoda y un compromiso con fecha. La parte difícil es la última.
\end{softbox}

\milcpaso{1}{milcMagenta}{\textbf{Las diez cosas (7 min).} Completa el inventario \textbf{con quién las hizo}. \emph{Y para las que quedaron en «no sé», haz algo esta semana:} \textbf{averigua una}. \emph{Cómo llegó el agua a tu barrio, quién consiguió el colegio, de cuándo es la vía, quién donó el lote del parque.} \textbf{Casi siempre hay una historia, y casi siempre involucra a gente que peleó por eso en reuniones aburridas} ---las del momento 7---.}
\milcpaso{2}{milcMagenta}{\textbf{Quién lo mantiene hoy (7 min).} Al frente de cada cosa, \emph{quién hace el trabajo ahora}. \textbf{Y marca las que quedaron en «nadie» o en «no sé».} \emph{Esas son, según Ostrom, las que están en camino de derrumbarse:} \textbf{lo que nadie sostiene no dura, aunque hoy se vea bien.}}
\milcpaso{3}{milcMagenta}{\textbf{Lo que está en peor estado (6 min).} Escoge \textbf{una} cosa de tu lista que esté peor que hace unos años ---el río, el parque, la biblioteca, una tradición, un espacio del colegio, una organización del barrio--- y escribe \emph{qué le pasó} y \emph{quién dejó de hacer el trabajo}. \textbf{Sin acusar a nadie:} \emph{casi siempre no es que alguien lo dañara, sino que quienes lo sostenían envejecieron y nadie los reemplazó.}}
\milcpaso{4}{milcMagenta}{\textbf{El compromiso con fecha (12 min).} Escribe \textbf{una devolución concreta} hacia algo que \textbf{vas a seguir usando}. \textbf{Cuatro condiciones:} \emph{que sea pequeña} ---esto no se hace con proyectos grandes---; \emph{que sea sostenible} ---algo que puedas repetir---; \emph{que tenga \textbf{fecha}}; y \emph{que sea hacia algo de tu lista, no hacia una causa lejana}. \textbf{Ejemplos reales:} \emph{ir a una reunión de la Junta, ayudar en la biblioteca, enseñarle a alguien algo que sabes, sumarte a una jornada del río, dejar en orden el espacio del colegio que usa el curso que viene, escribir la historia de algo que nadie ha escrito.} \emph{Y una advertencia:} \textbf{si tu compromiso es «cuidar el medio ambiente», no es un compromiso ---es un tema---.}}

\begin{drawbox}{milcMagenta}{Antes de cerrar tu inventario}
\begin{checklista}
\item Tus diez cosas incluyen \textbf{lo invisible}, no solo lo material.
\item Anotaste quién las hizo y \textbf{te propusiste averiguar una}.
\item Marcaste las que hoy \textbf{no las mantiene nadie}.
\item Identificaste una que está peor y \textbf{qué dejó de hacerse}.
\item Tu compromiso es \textbf{pequeño, repetible y con fecha}.
\item Es hacia algo que \textbf{vas a seguir usando}, no hacia una causa lejana.
\end{checklista}
\end{drawbox}

\begin{softbox}{milcMagenta}{milcGris}{Plantilla de trabajo}
\textbf{El inventario.} \emph{«Uso \_\_\_\_. Lo hizo \_\_\_\_. Hoy lo mantiene \_\_\_\_.»} \\[1mm]
\textbf{Lo que está en riesgo.} \emph{«\_\_\_\_ está peor porque dejó de \_\_\_\_.»} \\[1mm]
\textbf{Mi devolución.} \emph{«Voy a \_\_\_\_, el \_\_\_\_, y lo puedo repetir cada \_\_\_\_.»}
\end{softbox}

\begin{infoband}{milcMagenta}


\textbf{En tu cuaderno (Actividad 3 · CREA).} Título: «Actividad 3. Mi inventario de lo heredado». \textbf{Formato:} las diez cosas con sus dos columnas + la que está en peor estado y por qué + el compromiso con sus cuatro condiciones. \textbf{Extensión:} una página. \textbf{Sabes que terminaste cuando} tu compromiso \textbf{tiene fecha}.
\end{infoband}

\puente{El producto es un compromiso con fecha. \emph{Una deuda sin fecha es una manera elegante de no pagarla.}}

\milcestacion{milcMostaza}{Producto de la sesión}
\textbf{Inventario de lo heredado: diez cosas que usas y no hiciste, quién las sostiene hoy, y una devolución fechada}.
\begin{softbox}{milcMostaza}{milcGris}{Criterios de calidad}
(1) El inventario incluye bienes invisibles —derechos, lenguas, saberes— y no solo infraestructura. (2) Se identifica quién hizo cada cosa y quién la mantiene actualmente. (3) Se señalan las que hoy nadie sostiene, reconociéndolas como las que están en riesgo. (4) Se analiza un bien deteriorado sin acusar, identificando qué trabajo dejó de hacerse. (5) El compromiso es pequeño, repetible, fechado y dirigido a algo que la persona seguirá usando.
\end{softbox}

\puente{Antes de cerrar, revisa lo esencial: si tu devolución es hacia algo que no vas a volver a usar, es un favor y no una devolución.}

\milcestacion{milcVino}{Evaluación liberadora · cinco dimensiones}

\begin{softbox}{milcVino}{milcGris}{Sentido de la evaluación}
Evaluar no es solo revisar una nota. Es reconocer qué aprendí, qué cambió en mí, qué emociones manejé, cómo me relaciono con la ciudad y la comunidad, y qué le dejaré a quien viene detrás.
\end{softbox}

\textbf{Mi reflexión liberadora en cinco dimensiones.} Completa cada frase iniciada:

\small
\begin{tabularx}{\textwidth}{>{\bfseries\RaggedRight\arraybackslash}p{3.4cm}Y>{\RaggedRight\arraybackslash}p{4.6cm}}
\rowcolor{milcVino}\color{white}Dimensión & \color{white}\bfseries Mi respuesta (frame starter) & \color{white}\bfseries Algo que descubrí\\
1. Personal & Lo que más cambió en mí fue \linead{6.5cm} & \linead{4.2cm}\\
2. Emocional & La emoción más fuerte fue \linead{2.5cm} y la regulé \linead{3cm} & \linead{4.2cm}\\
3. Ciudadana & En democracia esto importa porque \linead{6cm} & \linead{4.2cm}\\
4. Local (Cartago) & En mi barrio sería útil para \linead{6.5cm} & \linead{4.2cm}\\
5. Intergeneracional & A mi abuela/abuelo le diría \linead{6.5cm} & \linead{4.2cm}\\
\end{tabularx}
\normalsize

\begin{infoband}{milcVino}
\textbf{Para la próxima sustentación.} Vuelve a esta tabla en seis meses. Verás que las respuestas cambian — no porque lo aprendido cambie, sino porque tú cambias. La evaluación liberadora es una conversación contigo a través del tiempo.
\end{infoband}


\puente{Cerramos con tres voces: el rostro que reclama un derecho, la abeja y el enjambre, y el cuidado de lo que se comparte.}

\milcestacion{milcGradeDark}{Triángulo de pensamiento · cierre filosófico}

\begin{softbox}{milcVino}{milcGris}{Dussel · lente del nosotros}
\textit{«El otro se revela realmente como otro\ldots como el pobre, el oprimido; el que, a la vera del camino, fuera del sistema, muestra su rostro sufriente y sin embargo desafiante: «¡Tengo hambre!, ¡tengo derecho a comer!»»}

{\footnotesize\itshape\color{milcNegro!70}--- Enrique Dussel · Filosofía de la liberación (1977)}

\emph{Hay una manera de leer este momento que sería falsa y conviene descartarla:} \textbf{«recibiste mucho, así que agradece».} \emph{Dussel apunta en otra dirección.} \textbf{Lo que uno recibió no llegó parejo:} \emph{no todos heredaron el mismo acueducto, la misma escuela ni las mismas oportunidades}, \textbf{y en este país hay quienes están \emph{fuera del sistema} de las cosas que otros dan por hechas.} \emph{Por eso la deuda de la que habla esta guía no es solo con el pasado ---con quienes hicieron lo que usas---:} \textbf{es también con quienes hoy no tienen acceso a eso.} \emph{Y la cita da el marco:} \textbf{lo que a ellos les falta no es tu generosidad ---es un derecho---.} \emph{Devolver, entonces, no es repartir lo que sobra: es sostener y ampliar lo que debería ser de todos.}

\textbf{De las diez cosas de mi lista, ¿cuáles no las tiene alguien que conozco?}
\end{softbox}
\linead{\linewidth}\\[1mm]
\linead{\linewidth}

\begin{softbox}{milcMostaza}{milcGris}{Estoicismo · Marco Aurelio · Meditaciones VI, 54 (c. 175 d.C.) · cuidado interior}
\textit{«Lo que no beneficia al enjambre, tampoco beneficia a la abeja.»}

{\footnotesize\itshape\color{milcNegro!70}--- Marco Aurelio · Meditaciones VI, 54 (c. 175 d.C.)}

\emph{Marco Aurelio la escribe en presente ---lo que le sirve al grupo hoy---.} \textbf{Ostrom le agrega el tiempo largo:} \emph{los bienes comunes que duraron generaciones tenían principios de diseño que los que se derrumbaron no tenían} (Ostrom, 1990). \textbf{Es decir: el enjambre no es solo el que está vivo ahora.} \emph{Incluye a los que hicieron el panal y a los que van a llegar.} \emph{Y esa es la forma más honesta de plantear la deuda a los diecisiete:} \textbf{no se le devuelve a quien te dio ---esa gente muchas veces ya no está---.} \emph{Se le devuelve al sistema, hacia adelante.} \textbf{Es la minga, otra vez, pero medida en generaciones:} \emph{la cuenta se lleva en turnos, y ahora empieza el tuyo.}

\textbf{Lo que voy a usar los próximos diez años, ¿quién lo va a estar sosteniendo?}
\end{softbox}
\linead{\linewidth}\\[1mm]
\linead{\linewidth}

\begin{softbox}{milcTurquesa}{milcGris}{Floridi · lente de la infoesfera}
\textit{«El florecimiento de las entidades informacionales y de toda la infoesfera debe promoverse preservando, cultivando y enriqueciendo sus propiedades.»}

{\footnotesize\itshape\color{milcNegro!70}--- Luciano Floridi · La ética de la información (2013; trad.)}

\textbf{Los tres verbos de Floridi son, exactamente, los tres trabajos que exige algo heredado:} \emph{\textbf{preservar} lo que ya existe, \textbf{cultivar} lo que necesita mantenimiento, \textbf{enriquecer} lo que se le puede agregar.} \emph{Y hay un bien común de tu generación que casi nadie cuenta como tal:} \textbf{lo que se sabe de un lugar.} \emph{La historia del barrio, cómo se consiguió el acueducto, quién fundó qué, las recetas, los nombres de las plantas, cómo se hacía un oficio.} \textbf{Eso vive en la cabeza de gente mayor y \emph{se pierde entero} cuando esa gente falta ---no queda ningún registro---.} \emph{Y ustedes tienen algo que ninguna generación anterior tuvo:} \textbf{la posibilidad de grabar una conversación de treinta minutos con un teléfono.} \emph{Preservar, para ustedes, es más barato que nunca, y por eso no hacerlo pesa más.}

\textbf{¿Qué se pierde de mi familia o mi barrio si nadie pregunta y nadie graba?}
\end{softbox}
\linead{\linewidth}\\[1mm]
\linead{\linewidth}

\begin{infoband}{milcNegro}
\textbf{Nota para el docente.} Las tres son citas textuales y llevan su referencia debajo. Puedes remitir a la obra en clase.
\end{infoband}

\milcestacion{milcVino}{Mi autoevaluación · marca una casilla por fila}

{\small
\setlength{\extrarowheight}{2.2pt}
\begin{tabularx}{\textwidth}{>{\RaggedRight\arraybackslash\bfseries}p{3.0cm}YYY}
\rowcolor{milcVino}\color{white}\bfseries Criterio & \color{white}\bfseries Lo logré & \color{white}\bfseries En proceso & \color{white}\bfseries Necesito apoyo\\[.6mm]
Comprensión del tema & \milcopcion{Explico las ideas con mis palabras.} & \milcopcion{Reconozco las ideas, falta precisión.} & \milcopcion{Necesito repasar lo básico.}\\[.8mm]
\rowcolor{milcGris}Calidad del inventario & \milcopcion{Puedo nombrar quién hizo y quién mantiene lo que uso, y tengo un compromiso concreto con fecha.} & \milcopcion{Reconozco que heredé cosas, pero mi compromiso es una intención sin fecha.} & \milcopcion{Sigo pensando que lo que hay simplemente estaba ahí.}\\[.8mm]
Aplicación y producto & \milcopcion{Mi producto cumple los criterios.} & \milcopcion{Está armado, con ajustes pendientes.} & \milcopcion{Aún está en construcción.}\\[.8mm]
\rowcolor{milcGris}Reflexión 5 dimensiones & \milcopcion{Respondí cada una con honestidad.} & \milcopcion{Respondí algunas con detalle.} & \milcopcion{Mis respuestas son muy generales.}\\[.8mm]
Triángulo de pensamiento & \milcopcion{Conecté las 3 voces con mi trabajo.} & \milcopcion{Conecté al menos una.} & \milcopcion{No las relaciono con mi proceso.}\\[.8mm]
\end{tabularx}}

\vspace{3mm}
\textbf{Hoy entendí que:}\\[1mm]
\linead{\linewidth}\\[1mm]
\linead{\linewidth}

\textbf{Mi compromiso para la próxima sesión es:}\\[1mm]
\linead{\linewidth}\\[1mm]
\linead{\linewidth}

\begin{softbox}{milcMostaza}{milcGris}{Cita de despedida · Marco Aurelio}
\textit{``El obstáculo en el camino se vuelve el camino. Lo que estorba al camino se vuelve camino.''}\\[1mm]
Cada error de hoy, bien observado, se convierte en pieza del camino que pisarás mañana.
\end{softbox}

\small
\begin{tabularx}{\textwidth}{>{\bfseries}p{3.6cm}Y}
\rowcolor{milcGradeDark!12}Nombre completo: & \linead{10cm}\\
Fecha: & \linead{4cm} \quad Curso/grupo: \linead{4cm}\\
Mi firma: & \linead{9cm}\\
\end{tabularx}
\normalsize

\begin{infoband}{milcGradeDark}
\textbf{Para la próxima sesión.} Dos cosas. \textbf{Primero, cumple la primera fecha de tu devolución} y trae \textbf{qué hiciste} ---no qué sentiste---. \emph{Y averigua una de las cosas que quedaron en «no sé»: cómo llegó el agua a tu barrio, quién consiguió el colegio, de cuándo es la vía.} \textbf{Segundo, y este año conviene hacerlo bien: pregunta en tu casa por lo que costó.} Averigua con alguien mayor \textbf{qué cosas de las que hoy usan no existían cuando llegaron} ---el agua, la luz, la escuela, la vía, el puesto de salud--- y \textbf{cómo se consiguieron}: \emph{quién peleó, cuánto tardó, si hubo que ir a Cali o a Bogotá, quiénes pusieron el trabajo}. \textbf{Y si puedes, grábalo.} \emph{Treinta minutos y un teléfono.} \textbf{Trae:} tu devolución cumplida y la historia. \emph{Vas a encontrar lo mismo que encontró Ostrom en cientos de casos, contado en primera persona: casi todo lo que hay lo consiguió un grupo pequeño de gente terca, casi nadie les agradeció, y casi ninguno de sus nombres está escrito en ninguna parte. La grabación que hagas puede ser el único registro que quede.}
\end{infoband}

\end{document}
